%%%%%%%%%%%%%%%%%%%%%%%%%%%%%%%%%%%%%%%%%%%%%%%%%%%%%%%%%%%%%%%%%%%%%%%%
%%%%%%%%%%%%%%%%%%%%%% Simple LaTeX CV Template %%%%%%%%%%%%%%%%%%%%%%%%
%%%%%%%%%%%%%%%%%%%%%%%%%%%%%%%%%%%%%%%%%%%%%%%%%%%%%%%%%%%%%%%%%%%%%%%%

%THIS FILE IS BASED ON THE TEMPLATE FROM "http://www.tedpavlic.com/post_resume_cv_latex_example.php"

%%%%%%%%%%%%%%%%%%%%%%%%%%%%%%%%%%%%%%%%%%%%%%%%%%%%%%%%%%%%%%%%%%%%%%%%
%% NOTE: If you find that it says                                     %%
%%                                                                    %%
%%                           1 of ??                                  %%
%%                                                                    %%
%% at the bottom of your first page, this means that the AUX file     %%
%% was not available when you ran LaTeX on this source. Simply RERUN  %%
%% LaTeX to get the ``??'' replaced with the number of the last page  %%
%% of the document. The AUX file will be generated on the first run   %%
%% of LaTeX and used on the second run to fill in all of the          %%
%% references.                                                        %%
%%%%%%%%%%%%%%%%%%%%%%%%%%%%%%%%%%%%%%%%%%%%%%%%%%%%%%%%%%%%%%%%%%%%%%%%

%%%%%%%%%%%%%%%%%%%%%%%%%%%% Document Setup %%%%%%%%%%%%%%%%%%%%%%%%%%%%

% Don't like 10pt? Try 11pt or 12pt
\documentclass[10pt]{article}
\usepackage{verbatim}

% The automated optical recognition software used to digitize resume
% information works best with fonts that do not have serifs. This
% command uses a sans serif font throughout. Uncomment both lines (or at
% least the second) to restore a Roman font (i.e., a font with serifs).
%\usepackage{times}
%\renewcommand{\familydefault}{\sfdefault}

% This is a helpful package that puts math inside length specifications
\usepackage{calc}

% Layout: Puts the section titles on left side of page
\reversemarginpar

%
%         PAPER SIZE, PAGE NUMBER, AND DOCUMENT LAYOUT NOTES:
%
% The next \usepackage line changes the layout for CV style section
% headings as marginal notes. It also sets up the paper size as either
% letter or A4. By default, letter was used. If A4 paper is desired,
% comment out the letterpaper lines and uncomment the a4paper lines.
%
% As you can see, the margin widths and section title widths can be
% easily adjusted.
%
% ALSO: Notice that the includefoot option can be commented OUT in order
% to put the PAGE NUMBER *IN* the bottom margin. This will make the
% effective text area larger.
%
% IF YOU WISH TO REMOVE THE ``of LASTPAGE'' next to each page number,
% see the note about the +LP and -LP lines below. Comment out the +LP
% and uncomment the -LP.
%
% IF YOU WISH TO REMOVE PAGE NUMBERS, be sure that the includefoot line
% is uncommented and ALSO uncomment the \pagestyle{empty} a few lines
% below.
%

%% Use these lines for letter-sized paper
\usepackage[paper=letterpaper,
            %includefoot, % Uncomment to put page number above margin
            marginparwidth=1.2in,     % Length of section titles
            marginparsep=.05in,       % Space between titles and text
            margin=1in,               % 1 inch margins
            includemp]{geometry}

%% Use these lines for A4-sized paper
%\usepackage[paper=a4paper,
%            %includefoot, % Uncomment to put page number above margin
%            marginparwidth=30.5mm,    % Length of section titles
%            marginparsep=1.5mm,       % Space between titles and text
%            margin=25mm,              % 25mm margins
%            includemp]{geometry}

%% More layout: Get rid of indenting throughout entire document
\setlength{\parindent}{0in}

\usepackage[shortlabels]{enumitem}

% Simpler bibsections for CV sections
% (thanks to natbib for inspiration)
%
% * For lists of references with hanging indents and no numbers:
%
%   \begin{bibsection}
%       \item ...
%   \end{bibsection}
%
% * For numbered lists of references (with hanging indents):
%
%   \begin{bibenum}
%       \item ...
%   \end{bibenum}
%
%   Note that bibenum numbers continuously throughout. To reset the
%   counter, use
%
%   \restartlist{bibenum}
%
%   at the place where you want the numbering to reset.

\makeatletter
\newlength{\bibhang}
\setlength{\bibhang}{1em}
\newlength{\bibsep}
 {\@listi \global\bibsep\itemsep \global\advance\bibsep by\parsep}
\newlist{bibsection}{itemize}{3}
\setlist[bibsection]{label=,leftmargin=\bibhang,%
        itemindent=-\bibhang,
        itemsep=\bibsep,parsep=\z@,partopsep=0pt,
        topsep=0pt}
\newlist{bibenum}{enumerate}{3}
\setlist[bibenum]{label=[\arabic*],resume,leftmargin={\bibhang+\widthof{[999]}},%
        itemindent=-\bibhang,
        itemsep=\bibsep,parsep=\z@,partopsep=0pt,
        topsep=0pt}
\let\oldendbibenum\endbibenum
\def\endbibenum{\oldendbibenum\vspace{-.6\baselineskip}}
\let\oldendbibsection\endbibsection
\def\endbibsection{\oldendbibsection\vspace{-.6\baselineskip}}
\makeatother

%% Reference the last page in the page number
%
% NOTE: comment the +LP line and uncomment the -LP line to have page
%       numbers without the ``of ##'' last page reference)
%
% NOTE: uncomment the \pagestyle{empty} line to get rid of all page
%       numbers (make sure includefoot is commented out above)
%
\usepackage{fancyhdr,lastpage}
\pagestyle{fancy}
\pagestyle{empty}      % Uncomment this to get rid of page numbers
\fancyhf{}\renewcommand{\headrulewidth}{0pt}
\fancyfootoffset{\marginparsep+\marginparwidth}
\newlength{\footpageshift}
\setlength{\footpageshift}
          {0.5\textwidth+0.5\marginparsep+0.5\marginparwidth-2in}
\lfoot{\hspace{\footpageshift}%
       \parbox{4in}{\, \hfill %
                    \arabic{page} of \protect\pageref*{LastPage} % +LP
%                    \arabic{page}                               % -LP
                    \hfill \,}}

% Finally, give us PDF bookmarks
\usepackage{color,hyperref}
\definecolor{darkblue}{rgb}{0.0,0.0,0.3}
\hypersetup{colorlinks,breaklinks,
            linkcolor=darkblue,urlcolor=darkblue,
            anchorcolor=darkblue,citecolor=darkblue}

\usepackage{url}


%%%%%%%%%%%%%%%%%%%%%%%% End Document Setup %%%%%%%%%%%%%%%%%%%%%%%%%%%%


%%%%%%%%%%%%%%%%%%%%%%%%%%% Helper Commands %%%%%%%%%%%%%%%%%%%%%%%%%%%%

%%% HEADING AT TOP OF CURRICULUM VITAE

% The title (name) with a horizontal rule under it
% (optional argument typesets an object right-justified across from name
%  as well)
%
% Usage: \makeheading{name}
%        OR
%        \makeheading[right_object]{name}
%
% Place at top of document. It should be the first thing.
% If ``right_object'' is provided in the square-braced optional
% argument, it will be right justified on the same line as ``name'' at
% the top of the CV. For example:
%
%       \makeheading[\emph{Curriculum vitae}]{Your Name}
%
% will put an emphasized ``Curriculum vitae'' at the top of the document
% as a title. Likewise, a picture could be included:
%
%   \makeheading[\includegraphics[height=1.5in]{my_picutre}]{Your Name}
%
% the picture will be flush right across from the name.
\newcommand{\makeheading}[2][]%
        {\hspace*{-\marginparwidth}%
         \begin{minipage}[t]{\textwidth+\marginparwidth+\marginparsep}%
             {\large \bfseries #2 \hfill #1}%
         \end{minipage}%
	\vspace{-1.3\baselineskip}}

%%% SECTION HEADINGS

% The section headings. Flush left in small caps down pseudo-margin.
%
% Usage: \section{section name}
\renewcommand{\section}[1]{%
    \pagebreak[3]%
    \vspace{1.3\baselineskip}%
    \phantomsection\addcontentsline{toc}{section}{#1}%
    \hspace{-\marginparwidth}\rule{\columnwidth+\marginparwidth}{1pt}\newline%
    \noindent\llap{\scshape\smash{\parbox[t]{\marginparwidth}{\hyphenpenalty=10000\raggedright #1}}}%
    \vspace{-\baselineskip}\par}

\renewcommand{\subsection}[1]{%
    \pagebreak[3]%
    \vspace{1.3\baselineskip}%
    \phantomsection\addcontentsline{toc}{section}{#1}%
    \noindent\llap{\scshape\smash{\parbox[t]{\marginparwidth}{\hyphenpenalty=10000\raggedright #1}}}%
    \vspace{-\baselineskip}\par}

%%% LISTS

% This macro alters a list by removing some of the space that follows the list
% (is used by lists below)
\newcommand*\fixendlist[1]{%
    \expandafter\let\csname preFixEndListend#1\expandafter\endcsname\csname end#1\endcsname
    \expandafter\def\csname end#1\endcsname{\csname preFixEndListend#1\endcsname\vspace{-0.6\baselineskip}}}

% These macros help ensure that items in outer-type lists do not get
% separated from the next line by a page break
% (they are used by lists below)
\let\originalItem\item
\newcommand*\fixouterlist[1]{%
    \expandafter\let\csname preFixOuterList#1\expandafter\endcsname\csname #1\endcsname
    \expandafter\def\csname #1\endcsname{\csname preFixOuterList#1\endcsname\let\oldItem\item\def\item{\pagebreak[2]\oldItem}}
    \expandafter\let\csname preFixOuterListend#1\expandafter\endcsname\csname end#1\endcsname
    \expandafter\def\csname end#1\endcsname{\let\item\oldItem\csname preFixOuterListend#1\endcsname}}
\newcommand*\fixinnerlist[1]{%
    \expandafter\let\csname preFixInnerList#1\expandafter\endcsname\csname #1\endcsname
    \expandafter\def\csname #1\endcsname{\let\oldItem\item\let\item\originalItem\csname preFixInnerList#1\endcsname}
    \expandafter\let\csname preFixInnerListend#1\expandafter\endcsname\csname end#1\endcsname
    \expandafter\def\csname end#1\endcsname{\csname preFixInnerListend#1\endcsname\let\item\oldItem}}

\newcommand{\sbullet}{\,\begin{picture}(1,1)(0,-3)\circle*{3}\end{picture}\ }

% An itemize-style list with lots of space between items
%
% Usage:
%   \begin{outerlist}
%       \item ...    % (or \item[] for no bullet)
%   \end{outerlist}
\newlist{outerlist}{itemize}{3}
    \setlist[outerlist]{label=\enskip\sbullet,leftmargin=*}
    \fixendlist{outerlist}
    \fixouterlist{outerlist}

% An environment IDENTICAL to outerlist that has better pre-list spacing
% when used as the first thing in a \section
%
% Usage:
%   \begin{lonelist}
%       \item ...    % (or \item[] for no bullet)
%   \end{lonelist}
\newlist{lonelist}{itemize}{3}
    \setlist[lonelist]{label=\enskip\sbullet,leftmargin=*,partopsep=0pt,topsep=0pt}
    \fixendlist{lonelist}
    \fixouterlist{lonelist}

% An itemize-style list with little space between items
%
% Usage:
%   \begin{innerlist}
%       \item ...    % (or \item[] for no bullet)
%   \end{innerlist}
\newlist{innerlist}{itemize}{3}
    \setlist[innerlist]{label=\enskip\sbullet,leftmargin=*,parsep=0pt,itemsep=0pt,topsep=0pt,partopsep=0pt}
    \fixinnerlist{innerlist}

% An environment IDENTICAL to innerlist that has better pre-list spacing
% when used as the first thing in a \section
%
% Usage:
%   \begin{loneinnerlist}
%       \item ...    % (or \item[] for no bullet)
%   \end{loneinnerlist}
\newlist{loneinnerlist}{itemize}{3}
    \setlist[loneinnerlist]{label=\enskip\sbullet,leftmargin=*,parsep=0pt,itemsep=0pt,topsep=0pt,partopsep=0pt}
    \fixendlist{loneinnerlist}
    \fixinnerlist{loneinnerlist}

%%% EXTRA SPACE

% To add some paragraph space between lines.
% This also tells LaTeX to preferably break a page on one of these gaps
% if there is a needed pagebreak nearby.
\newcommand{\blankline}{\quad\pagebreak[3]}
\newcommand{\halfblankline}{\quad\vspace{-0.9\baselineskip}\pagebreak[3]}

%%%%%%%%%%%%%%%%%%%%%%%% End Helper Commands %%%%%%%%%%%%%%%%%%%%%%%%%%%

%%%%%%%%%%%%%%%%%%%%%%%%% Begin CV Document %%%%%%%%%%%%%%%%%%%%%%%%%%%%

\begin{document}
\makeheading[\emph{Curriculum vitae}]{Connor R. Schenck}

\section{Office \newline Address}

% NOTE: Mind where the & separators and \\ breaks are in the following
%       table. Table is one row made up of three parboxes. The left
%       parbox has address info, the middle parbox has a vertical bar,
%       and the right parbox has phone and electronic contact
%       information.
%
% MACROS: \rcollength is the width of the right column of the table
%             (adjust it to your liking; default is 1.85in).
%         \spacewidth is width of area between left and right boxes.
%         \spacechar is character used to produce perforated vertical
%             boundary between boxes.
%
\newlength{\rcollength}\setlength{\rcollength}{2.8in}%
\newlength{\spacewidth}\setlength{\spacewidth}{20pt}
\newcommand\spacechar{$|$}
%
\vspace{1mm}
\begin{tabular}[t]{@{}p{\textwidth-\rcollength-\spacewidth}@{}p{\spacewidth}@{}p{\rcollength}}%

% Address box
\parbox{\textwidth-\rcollength-\spacewidth}{%
Computer Science and Engineering\\
University of Washington \\
185 Stevens Way \\
Seattle WA 98195-2350}

& \parbox{\spacewidth}{} &

% Non-snail-mail contact information
\parbox{\rcollength}{%
\begin{tabular}{r l}
\textit{Phone}:& \hspace{-0.10in}(515) 257-6335 \\
\textit{E-mail}:& \hspace{-0.10in}schenckc@cs.washington.edu\\
\textit{URL}:& \hspace{-0.10in}{\small\url{homes.cs.washington.edu/~schenckc}}\\
\end{tabular}}


\end{tabular}

\section{Research Interests}

\textbf{Artificial Intelligence with an emphasis on learning and adaptation} including: Machine Learning, Robotics, Developmental Robotics, Computational Perception, Robotic Manipulation, Autonomous Systems, Humanoid Robotics, and Artificial Intelligence for Video Games.


\section{Education}

\textbf{University of Washington} \hfill Seattle, WA

\begin{outerlist}

\item[] \emph{Doctor of Philosophy} \hfill {Summer 2018}
        \begin{innerlist}
	\item Major: Computer Science \& Engineering
        \item Thesis Field: Robotics
        \item Advisor: Dr. Dieter Fox
        \end{innerlist}

\end{outerlist}

\blankline

\textbf{Iowa State University of Science and Technology} \hfill Ames, IA
\begin{outerlist}

\item[] \emph{Master of Science} \hfill {Summer 2013}
        \begin{innerlist}
	\item Major: Computer Science
	\item Co-Major: Human-Computer Interaction
        \item Thesis Field: Developmental Robotics
        \item Advisor: Dr. Alexander Stoytchev
        \end{innerlist}

\halfblankline

\item[] \emph{Bachelor of Science} \hfill {Fall 2011}
        \begin{innerlist}
	\item Major: Computer Science
        \item Emphasis: Artificial Intelligence and Machine Learning
        \item GPA: 3.92/4.0 \emph{Summa Cum Laude}
        \end{innerlist}

\end{outerlist}

\section{Honors \& Awards}

\begin{innerlist}
	\item National Science Foundation Graduate Research Fellow
	\item Goldwater Scholar
	\item President's Leadership Scholarship
	\item President's Award for Competitive Excellence
	\item H Stuart Kuyper Engineering Endowment Fund
	\item Dean's List
	\item Highest 2\% ISU Engineering Freshmen
\end{innerlist}

\section{Publications \& Presentations}

\begin{outerlist}

\item[] \textbf{PhD Thesis}
	\begin{innerlist}
		\item \textbf{Schenck, C.}, ``Liquids \& Robots: An Investigation of Techniques for Robotic Interaction with Liquids,'' PhD Thesis, Paul G. Allen School of Computer Science \& Engineering, University of Washington, Seattle, WA, August 2018.
	\end{innerlist}

\halfblankline

\item[] \textbf{Master's Thesis}
	\begin{innerlist}
		\item \textbf{Schenck, C.}, ``Intelligence Tests for Robots: Solving Perceptual Reasoning Tasks with a Humanoid Robot,'' M.S. Thesis, Department of Computer Science \& Human-Computer Interaction Program, Iowa State University, Ames, IA, July 2013.
	\end{innerlist}

\halfblankline

\item[] \textbf{Refereed Journal Articles}
	\begin{innerlist}
	    \item \textbf{Schenck, C.}, and Fox, D, ``Perceiving and Reasoning About Liquids Using Fully-Convolutional Networks,'' \emph{International Journal of Robotics Research (IJRR)}, Vol.~37, No.~4--5, pp.~452--471, April, 2018.
		\item \textbf{Schenck, C.}, Sinapov, J., Johnston, J., and Stoytchev A., ``Which Object Fits Best? Solving Matrix Completion Tasks with a Humanoid Robot,'' \emph{IEEE Transactions on Autonomous Mental Development}, Vol.~6, No.~3, pp.~226--240, 2014.
		\item Sinapov, J., \textbf{Schenck, C.}, Staley, K., Sukhoy, V., and Stoytchev A., ``Grounding Semantic Categories in Behavioral Interactions: Experiments with 100 Objects,'' \emph{Journal of Robotics and Autonomous Systems}, Vol.~62, No.~5, pp.~632--645, 2014.	
		\item \textbf{Schenck, C.}, Sinapov, J., and Stoytchev, A., ``Which Object Comes Next? Grounded Order Completion by a Humanoid Robot,'' \emph{Journal of Cybernetics and Information Technologies}, Vol.~12, No.~3, pp. 5--16, 2012.
		\item Sinapov, J., Bergquist, T., \textbf{Schenck, C.}, Ohiri, U., Griffith, S., and Stoytchev, A., ``Interactive Object Recognition Using Proprioceptive and Auditory Feedback,'' \emph{International Journal of Robotics Research}, Vol. 30, No. 10, pp. 1250--1262, 2011.
	\end{innerlist}

\halfblankline

\item[] \textbf{Conference Articles (peer reviewed)}
	\begin{innerlist}
	\item \textbf{Schenck, C.}, and Fox, D., ``SPNets: Differentiable Fluid Dynamics for Deep Neural Networks,'' In \emph{Proceedings of the Second Conference on Robot Learning (CoRL)}, Zurich, Switzerland, October 29--31, 2018.
	\item \textbf{Schenck, C.}, Tompson, J., Fox, D., and Levine, S., ``Learning Robotic Manipulation of Granular Media,'' In \emph{Proceedings of the First Conference on Robot Learning (CoRL)}, Mountain View, CA, USA, November 13--15, 2017.
	\item \textbf{Schenck, C.}, and Fox, D, ``Reasoning About Liquids via Closed-Loop Simulation,'' In \emph{Robotics: Science \& Systems (RSS)}, Cambridge, MA, USA, July 12--16, 2017.
	\item \textbf{Schenck, C.}, and Fox, D, ``Visual Closed-Loop Control for Pouring Liquids,'' In \emph{Proceedings of the International Conference on Experimental Robotics (ICRA)}, Singapore, May 29 -- June 3, 2017.
	\item \textbf{Schenck, C.}, and Fox, D, ``Towards Learning to Perceive and Reason About Liquids,'' In \emph{Proceedings of the International Symposium on Experimental Robotics (ISER)}, Tokyo, Japan, October 3--6, 2016.
	\item Sinapov, J., \textbf{Schenck, C.}, and Stoytchev, A, ``Learning Relational Object Categories Using Behavioral Exploration and Multimodal Perception,'' In \emph{Proceedings of the 2014 IEEE International Conference on Robotics and Automation (ICRA)}, Hong Kong, China, May 31 -- June 7, 2014.
	\end{innerlist}

\halfblankline

\item[] \textbf{Workshop Presentations \& Articles (peer reviewed)}
	\begin{innerlist}
	\item \textbf{Schenck, C.}, and Fox, D., ``SPNets: Modeling Position Based Fluids using Smooth Particle Networks,'' In \emph{Proceedings of Robotics Science \& Systems (RSS) 2018 Workshop Learning and Inference in Robotics: Integrating Structure, Priors and Models}, Pittsburgh, PA, USA, June 29, 2018.
	\item \textbf{Schenck, C.}, and Fox, D., ``Detection and Tracking of Liquids with Fully Convolutional Networks,'' In \emph{Proceedings of Robotics Science \& Systems (RSS) 2016 Workshop Are the Skeptics Right? Limits and Potentials of Deep Learning in Robotics}, Ann Arbor, Michigan, USA, June 18, 2016.
	\item \textbf{Schenck, C.}, and Sinapov, A., ``The Object Pairing and Matching Task: Toward Montessori Tests for Robots,'' In \emph{Proceedings of the Humanoids 2012 Workshop on Developmental Robotics}, Okaka, Japan, November 29, 2012.
	\item \textbf{Schenck, C.}, Sinapov, J., and Stoytchev, A., ``Which Object Comes Next? Grounded Order Completion by a Humanoid Robot,'' \emph{AIMSA Workshop: Advances in Robot Learning and Human-Robot Interaction}, Varna, Bulgaria, September 12, 2012.
	\item Bergquist, T., \textbf{Schenck, C.}, Ohiri, U., Sinapov, J., Griffith, S., and Stoytchev, A., ``Interactive Object Recognition Using Proprioceptive Feedback,'' In \emph{Proceedings of the IROS 2009 Workshop: Semantic Perception for Mobile Manipulation}, St.~Louis, MO, October 15, 2009.
	\end{innerlist}

\halfblankline

\item[] \textbf{Symposium Presentations \& Invited Talks}
	\begin{innerlist}
	    \item \textbf{Schenck, C.}, ``Towards Perceiving and Manipulating Liquids with Real Robots,'' \emph{Invited Talk}, University of California at Berkeley, Berkeley, California, April 21, 2017.
		\item \textbf{Schenck, C.}, ``Object Recognition Using Proprioceptive and Auditory Feedback,'' \emph{Iowa State University Undergraduate Research Symposium}, April 19, 2011.
		\item Bergquist, T., \textbf{Schenck, C.}, and Ohiri, U., ``Interactive Object Recognition Using Proprioceptive and Auditory Feedback,'' \emph{Iowa State University REU Research Symposium}, Ames, IA, July 31, 2009.
	\end{innerlist}

\end{outerlist}

\section{Research Experience}

\textbf{Nvidia Corporation} \hfill Seattle, WA
\vspace{-0.17cm}
\begin{outerlist}

\item[] \emph{Research Intern} \hfill Winter 2018 (in progress)\\
	\emph{Robotics Research Lab}\vspace{-0.1cm}
        \begin{innerlist}
		\item Conducted research in robotics on combining deep learning with fluid simulation.
		\item Developed new particle convolution layer for performing convolutions on unordered particle sets.
		\item Constructed a fluid simulator using only deep neural networks.
        \end{innerlist}

\end{outerlist}

\blankline

\textbf{Google, Inc.} \hfill Mountain View, CA
\vspace{-0.17cm}
\begin{outerlist}

\item[] \emph{Software Engineering Intern} \hfill Winter/Spring 2017\\
	\emph{Google Brain}\vspace{-0.1cm}
        \begin{innerlist}
		\item Conducted research in robotics focusing on applying deep learning to robotic tasks.
		\item Published the results of the research in a peer-reviewed academic conference.
        \end{innerlist}

\item[] \emph{Software Engineering Intern} \hfill Summer 2015\\
	\emph{Google X}\vspace{-0.1cm}
        \begin{innerlist}
		\item Created a tool to help analyze the performance of various modules of the self-driving car's codebase.
		\item Applied machine learning to detect yielding of other cars to the self-driving car.
        \end{innerlist}

\end{outerlist}

\blankline

\textbf{Iowa State University} \hfill Ames, IA 
\vspace{-0.17cm}
\begin{outerlist}
\item[] \emph{Software Developer} \hfill Summer 2011 to Spring 2012\\
	\emph{Department of Statistics}\vspace{-0.1cm}
        \begin{innerlist}
		\item Implemented mutli-byte character support in a text coding application
		\item Debugged the large code-base collaboratively made with many other developers and released the application
        \end{innerlist}

\item[] \emph{Research Assistant} \hfill Summer/Fall 2010\\
	\emph{Virtual Reality Applications Center}
        \begin{innerlist}
		\item Developed a tool tracking application for factory workers
		\item Worked with factory workers to incorporate feedback
        \end{innerlist}

\item[] \emph{Research Experience for Undergraduates} \hfill Summer 2009\\
	\emph{Department of Human-Computer Interaction}
        \begin{innerlist}
		\item Chose a research topic and created a research question
		\item Conducted research in Developmental Robotics
		\item Coordinated tasks between group members
		\item Wrote and presented a report on the research 
        \end{innerlist}

\end{outerlist}



\section{Teaching Experience}

\textbf{Teaching Assistant}

\textbf{University of Washington} \hfill Seattle, WA

\begin{outerlist}

\item[] \emph{CSE 455: Computer Vision} \hfill Spring 2018 
        \begin{innerlist}
	\item Aided students from a wide variety of backgrounds in learning concepts in computer vision.
	\item Provided feedback to students via in class activities, assignments, and tests.
        \end{innerlist}

\item[] \emph{CSE 415: Introduction to AI for Non-Majors} \hfill Fall 2017 
        \begin{innerlist}
	\item Aided students from a wide variety of backgrounds in learning concepts in AI.
	\item Provided feedback to students via in class activities, assignments, and tests.
        \end{innerlist}

\end{outerlist}

\blankline

\textbf{Iowa State University} \hfill Ames, IA

\begin{outerlist}

\item[] \emph{ComS/HCI 575x: Computational Perception} \hfill Spring 2010 
        \begin{innerlist}
	\item Assisted students in understanding problems in Computer Vision through weekly office 
 hours and appointments
        \end{innerlist}

\halfblankline

\item[] \emph{CprE 185: Introduction to Programming in C} \hfill Fall 2009 
        \begin{innerlist}
	\item Designed lessons and projects for students in order to enhance understanding of concepts
	\item Aided students in understanding the basics of C programming
        \end{innerlist}

\halfblankline

\item[] \emph{ComS 227: Introduction to Object-Oriented Programming} \hfill Spring 2009 
        \begin{innerlist}
	\item Aided students in understanding concepts of object-oriented programming
	\item Proctored and graded exams, quizes, and homeworks
        \end{innerlist}

\end{outerlist}

\section{Technical Skills}

\begin{outerlist}

\item[] \textbf{Programming Languages, Libraries and Software Applications}
        \begin{innerlist}
		\item C, C++, Java, Matlab, Cuda, Unix shell scripting, \LaTeX, Linux, Windows, GNU make, OpenCV, Matlab Image Processing Toolkit, Java Swing, Microsoft Foundation Classes, Robot Operating System, Weka, Scikit-Learn, PyTorch, Caffe, Tensorflow
        \end{innerlist}

\halfblankline

\item[] \textbf{Technical Training \& Proficiency}
        \begin{innerlist}
		\item Human-Computer Interaction, Artificial Intelligence, Machine Learning, Computer Vision, Computational Perception, Algorithms, Probability \& Statistics, Computational Randomness, Probabilistic Robotics, Robotic Manipulation and Planning
        \end{innerlist}

\end{outerlist}

\section{Leadership}

\begin{outerlist}

\item[] \emph{Iowa State University Game Design Competition} \hfill Fall 2010/Spring 2011
        \begin{innerlist}
		\item Lead a team of 11 developers, artists, and business majors to design and implement a video game
		\item Won first place (\$10,000) in the PC and Console category
        \end{innerlist}

\halfblankline

\item[] \emph{Presiden's Leadership Class} \hfill Fall 2008
        \begin{innerlist}
		\item Group of 30 freshmen selected to meet with President Geoffroy of Iowa State University Once per week
		\item Developed leadership skills
        \end{innerlist}

\end{outerlist}

\begin{comment}
\section{Notes for Jim}

\begin{outerlist}

	\item I don't really like the last section here, but I want that information to stay, is there a good way to put that information somewhere else in the document?
	\item Multiple sections (Teaching Experience, Research Experience, and Leadership) have bullet points after entry that explain what I did (e.g., after ComS/HCI 575x it says ``Assisted students in understanding...''). When I was looking at other people's cv's, some people did that and some people did not. Should I have those bullet points or get rid of all/some of them?
	\item In presentations, the second one from the bottom was Taylor's presentation in St. Louis. I helped him prepare it and I was one of the main authors on the paper, but should I put that down since I didn't actually present it? Does ``Presentations'' mean presentations that I've given or times when my work has been presented?
	\item The Honors \& Awards section feels weird being as close to the top as it is, but this is the only way in which I could think of to include my NSF fellowship near the top. Is there a better way to do this? Maybe I should put it under research experience?
	\item One final note, the second listed publication and the first presentation are both dated in the future, but since they will have happened by the time I send this I wanted to include them so I didn't forget.

\end{outerlist}
\end{comment}


\end{document}

%%%%%%%%%%%%%%%%%%%%%%%%%% End CV Document %%%%%%%%%%%%%%%%%%%%%%%%%%%%%

%----------------------------------------------------------------------%
% The following is copyright and licensing information for
% redistribution of this LaTeX source code; it also includes a liability
% statement. If this source code is not being redistributed to others,
% it may be omitted. It has no effect on the function of the above code.
%----------------------------------------------------------------------%
% Copyright (c) 2007, 2008, 2009, 2010, 2011 by Theodore P. Pavlic,
% modified by Connor R. Schenck 2012
%
% Unless otherwise expressly stated, this work is licensed under the
% Creative Commons Attribution-Noncommercial 3.0 United States License. To
% view a copy of this license, visit
% http://creativecommons.org/licenses/by-nc/3.0/us/ or send a letter to
% Creative Commons, 171 Second Street, Suite 300, San Francisco,
% California, 94105, USA.
%
% THE SOFTWARE IS PROVIDED "AS IS", WITHOUT WARRANTY OF ANY KIND, EXPRESS
% OR IMPLIED, INCLUDING BUT NOT LIMITED TO THE WARRANTIES OF
% MERCHANTABILITY, FITNESS FOR A PARTICULAR PURPOSE AND NONINFRINGEMENT.
% IN NO EVENT SHALL THE AUTHORS OR COPYRIGHT HOLDERS BE LIABLE FOR ANY
% CLAIM, DAMAGES OR OTHER LIABILITY, WHETHER IN AN ACTION OF CONTRACT,
% TORT OR OTHERWISE, ARISING FROM, OUT OF OR IN CONNECTION WITH THE
% SOFTWARE OR THE USE OR OTHER DEALINGS IN THE SOFTWARE.
%----------------------------------------------------------------------%
